\chapter{Complex Amplitudes, and Where the Coherence Is}
\label{sec:quantum}
% ===================================================================

\begin{remark}[Status of this chapter]
\label{rmk:ch09status}
This chapter is the least finished thing in this part of the book, and the
reader should treat it accordingly.  An earlier version presented the
complex instance of the decoration as delivering quantum dynamics
outright.  It does not, and the reason it does not is structural rather
than technical --- it will not be repaired by working harder at the same
construction.  What has changed since that version is that the structural
reason is now known precisely, that it turns out to have two independent
sources rather than one, and that locating it exposes a place the
coherence \emph{can} come from which is not among the routes the chapter
previously offered.  That place is the equational theory.  This is
progress and it is not a result: what follows states a construction,
states an obstruction, and states where the construction puts what the
obstruction takes away.
\end{remark}

\section{The complex instance}

Take $\Semi = \CC$ in Definition~\ref{def:decoration}.  The amplitude of a
path is the product of its step amplitudes and the transition amplitude
from $P$ to $Q$ is the sum over paths,
\[
  \langle Q \mid P \rangle \;=\; \sum_{\gamma : P \to Q} W(\gamma),
\]
with $|\langle Q \mid P\rangle|^2$ the candidate transition probability.
This much is a construction and is not in doubt.  It is also, so far, only
bookkeeping: complex numbers have been attached to steps and multiplied
along paths.

What would make it quantum is \emph{cancellation}.  Two paths $\gamma_1,
\gamma_2$ from $P$ to $Q$ with $W(\gamma_1) = -W(\gamma_2)$ should
annihilate, so that the transition has zero probability despite having two
distinct causal histories.  That is the phenomenon; everything else is
notation.

\section{Two obstructions, of different kinds}
\label{sec:obstruction}

Cancellation \emph{between paths} is what this framework, as built,
forbids.  There are two reasons and they are worth keeping apart, because
one is about the semantics we chose and the other is about the
construction we built.

\subsection{The semantic obstruction}

\begin{remark}[Interference versus branching time]
\label{rmk:obstruction}
The equivalence on which this entire development rests is bisimulation,
and bisimulation is a \emph{branching-time} notion: it distinguishes
systems by the shape of their tree of possibilities, not merely by the set
of runs they admit.  Two distinct paths to the same term are, for
bisimulation, two distinct pieces of structure, and no assignment of
weights to steps can make a branching-time equivalence identify them ---
still less make them annul one another.  Interference adds and
\emph{cancels} contributions to a shared outcome, which is a statement
about runs; bisimulation refuses to pool runs in the first
place~\cite[\S13]{CE}.
\end{remark}

The obstruction is worth stating sharply because it is easy to disguise.
One can write down $\sum_\gamma W(\gamma)$ and observe that terms cancel;
this is arithmetic and always available.  What one cannot do, without
changing the semantics, is claim that the cancellation is \emph{observed}
by the theory --- that the resulting system is behaviourally
indistinguishable from one in which the transition never occurs.  Under
bisimulation it is not.  The two paths remain visible.

\subsection{The constructive obstruction}

The second reason is independent of the choice of equivalence, and it was
not available when this chapter was first written.  It comes from asking
what the complex instance has to do in order to be a lift of the real one.

Chapter~\ref{ch:weight} builds, for $\Semi = \CC$, one jump operator per
rule and refinement class.  Its coefficients have to be fixed somehow, and
the natural-looking choice --- assign an amplitude to each derivation and
sum the amplitudes over derivations reaching a common target --- has a
consequence that is easy to miss.  If $m$ distinct redexes of one class
carry a configuration to the same successor, summing amplitudes and
squaring at the end gives that channel weight $m^2\lambda$, where the real
instance gives it $m\lambda$.  The complex construction would then fail to
agree with the real one on a quantity both of them compute.

\begin{definition}[Conservativity]
\label{def:conservativity}
A complex weighting is \emph{conservative} over the corresponding real one
when, at zero Hamiltonian, the populations it evolves are exactly those
the real weighting evolves.
\end{definition}

Conservativity forces the other choice: sum the \emph{rates} over
derivations and take one square root, which is
Definition~\ref{wt:def:jump}.  And that choice removes cross-derivation
coherence by construction.  The two paths do not cancel --- not because a
branching-time equivalence declines to pool them, but because the only
normalisation under which the complex reading lifts the classical one adds
them underneath the square root, where nothing can cancel.

\begin{remark}[What the two obstructions do and do not share]
\label{rmk:two-obstructions}
They agree in their verdict and disagree in everything else.
Remark~\ref{rmk:obstruction} is a claim about \emph{observation}: the
cancellation could be computed but would not be seen.  The constructive
obstruction is a claim about \emph{arithmetic}: under the only
normalisation that lifts the classical semantics, there is no cancellation
to compute.  The second is the stronger, and it is worth noticing that it
survives a change of equivalence.  A reader who takes the first route of
Section~\ref{sec:whatwouldchange} --- move to linear time --- has removed
the semantic obstruction and still has this one.  The principle producing
it is conservativity, not physics, and a construction declining to be
conservative over its own classical instance would be buying interference
by giving up the claim that these are two instances of one thing.
\end{remark}

\begin{remark}[What was previously claimed about density matrices]
An earlier version of this chapter proposed letting the vectorial account
$\vect{A}$ become a positive semidefinite matrix $\rho$, with
$\mathrm{tr}(\rho)=1$ read as the quantum analogue of resource
conservation.  This should not be retained.  A resource account is an
element of an ordered monoid, tracking what may be spent; a density matrix
is a state, tracking what may be observed.  They are not the same kind of
thing, and identifying them conflates the conservation law of
Proposition~\ref{prop:firstlaw} --- which is about value under conversion
--- with the normalisation of a probability distribution.  Whatever the
right quantum statement is, it is not that one.
\end{remark}

\section{Where the coherence is}
\label{sec:coherence}

Both obstructions are about coherence \emph{between rewrites}.  Neither
says anything about coherence from any other source, and a presentation
has another source.

A GSLT is a triple: a signature, a set of equations, and a graph of
rewrite rules.  The rewrites are directed and costed, so they are
naturally dissipative.  The equations are symmetric and cost-free, so they
are naturally unitary.  A quantum reading has exactly two slots to fill,
and the presentation hands over exactly two kinds of relation to fill them
with.  Definition~\ref{wt:def:hamiltonian} does the filling: the rewrites
supply the jump operators, and equations \emph{withheld from the quotient}
supply a Hermitian hopping term.

The qualification carries the whole weight.  An equation quotiented into
the state space contributes nothing, because the two configurations it
relates have become one.  Most presentations quotient everything.  So most
weighted theories have $H = 0$, and by
Proposition~\ref{wt:prop:interference} that is exactly the case in which
the complex instance is the real instance wearing different notation.
Principle~\ref{wt:prin:slogan} states the consequence: a weighted theory
is quantum exactly to the extent that its presentation withholds equations
from the quotient.

Three things follow, in increasing order of how much they change what this
book has already said.

\paragraph{Complexification is not free, and now one can see what it costs.}
The earlier version of this chapter treated the complex instance as a
choice of semiring and nothing more, so that a modeller who wanted quantum
dynamics had only to select $\CC$.  That is not so.  Selecting $\CC$ over a
fully quotiented presentation changes nothing whatever.  What has to change
is the \emph{presentation} --- one must decline to quotient something, and
say with what amplitude it hops --- and that is a modelling decision made
in the object language, visible in the source text of a theory, rather
than a choice of coefficient ring made about it.

\paragraph{The obstruction has been relocated, not removed.}
Nothing above recovers cancellation between rewrites, and nothing above
should be read as claiming that a coherent hopping term makes a rewrite
theory quantum in the full sense.  What it does is identify the one place
in a presentation where a complex weighting has purchase, and thereby
convert ``we cannot get interference'' into the sharper and more useful
``interference lives in the equations and not in the rewrites, and here is
what putting it there costs''.  The honest summary is that the framework
is quantum in its equations and classical in its rewrites.  Whether that
division is a defect or a discovery is not something this chapter can
settle.

\paragraph{The ontology of Part~\ref{part:sciencebot} survives.}
This is the largest change, and it is a change to what the book concedes
rather than to what it claims.  The first route below --- move to a
linear-time semantics --- was previously the only route recovering
interference honestly, and its cost was stated plainly: the ontology in
which bisimulation classes are the real things does not survive the move.
That concession was live and nothing retracted it.  It can now be
retracted, because coherence sited in the equational theory never asks
bisimulation to pool runs.  Two configurations related by a withheld
equation are two configurations, distinguished exactly as branching time
distinguishes them, and the hopping between them is additional structure
on the state space rather than an identification of paths through it.  One
may have a nonzero Hamiltonian and keep the ontology.

\begin{remark}[What is still not answered]
\label{rmk:coherence-gaps}
Two things, and they should not be run together with the above.  First,
nothing here says which equations a physical theory would withhold, or
where the hopping amplitudes come from; the construction says where the
slot is and leaves it to the modeller, exactly as the choice of semiring
is left to the modeller.  Second, the notion of equivalence appropriate to
a theory with nonzero Hamiltonian is not known.  Rate bisimulation is the
obvious candidate at $H = 0$, and by
Proposition~\ref{wt:prop:interference} it transfers there without further
work; when $H \neq 0$ a relation on configurations must also respect
coherences, and no candidate is on offer.  That is the precise form in
which the branching-time question returns, and it is a better form than
the one this chapter started with.
\end{remark}

\section{What would have to change}
\label{sec:whatwouldchange}

For cancellation \emph{between rewrites}, three routes remain open and we
have not chosen among them.  The fourth item is not one of them; it is
where Section~\ref{sec:coherence} has already gone, recorded here so that
the list is not read as exhausting the options.

\begin{enumerate}[itemsep=5pt]
  \item \textbf{Move to linear time.}  Replace bisimulation by a
        trace-based equivalence and accept the loss of the branching-time
        distinctions the rest of this book depends on.  This removes the
        semantic obstruction and leaves the constructive one
        (Remark~\ref{rmk:two-obstructions}), so it is a smaller step than
        it once looked and buys less: one would additionally have to
        abandon conservativity, which is to say abandon the claim that the
        real and complex readings are two instances of one construction.

  \item \textbf{Quotient late.}  Keep bisimulation as the semantics and
        introduce cancellation as a separate extraction laid on top ---
        the Born rule as an additional map out of the model rather than a
        property of it.  This is the categorical-quantum-mechanics
        pattern, where the interaction structure yields a semiring of
        scalars and the probability rule is applied
        afterwards~\cite{abramsky}.  It is the most conservative route and
        the least explanatory.

  \item \textbf{Attach amplitudes to derivations.}  Cancellation between
        rewrites requires that two routes to a common contractum be able
        to carry different phases, which means a decoration indexed by
        derivations rather than by refinement classes.  That is a strictly
        larger construction and not a repair of this one; whether it can
        be given without destroying the locality that makes the present
        construction executable, and without making the decoration grow
        with the term, is the open question.

  \item \textbf{Withhold equations from the quotient} ---
        Section~\ref{sec:coherence}.  This is not a route to cancellation
        between rewrites and does not claim to be.  It is the observation
        that a presentation has a second place to put coherence, that this
        place is untouched by either obstruction, and that using it costs
        nothing the book has already spent.
\end{enumerate}

\section{Simulation, in the meantime}

The stochastic instance is unobstructed, and it is what the implementation
actually runs.  The classical Gillespie algorithm~\cite{gillespie1977exact}
simulates a continuous-time Markov chain by computing exit rates from the
current state, sampling a waiting time, and selecting the next rule
proportionally to rate.  With $\Semi = \Real_{\geq 0}$ and a dynamic
decoration (Definition~\ref{def:dyndec}) supplying the rates, this is
precisely a decorated GSLT stepped forward, with the qualification of
Theorem~\ref{wt:thm:markov}: the chain is over configurations and not over
terms, and the simulator advances the term, updates the table, debits the
account and records the trace as one step of one process.

The quantum case is no longer a loop in search of a semantics, and the
earlier version of this chapter was right to refuse it and wrong about
what the refusal was waiting for.  It was waiting for a fixed generator.
Quantum analogues of the Gillespie algorithm exist for open systems
described by Lindblad master equations~\cite{cao2021quantum}, and such an
analogue requires an operator assembled once and exponentiated; a
decoration whose table changes as the system runs does not obviously
supply one.  Definition~\ref{wt:def:hilbert} supplies it, by indexing the
basis with configurations rather than with terms, so that a table update
is an edge of the generator rather than a modification of it.  With that
in place the unravelling is well posed, and
Theorem~\ref{wt:thm:degeneration} says the classical algorithm is its
degeneration at zero Hamiltonian --- so the two simulators are one
simulator, and ``Gillespie-inspired'' is a statement about a limit rather
than a resemblance.

What remains withheld is the interpretation and not the computation.  One
may run the unravelling, and the numbers coming out will be the numbers of
a quantum dynamical semigroup.  Whether the system they describe is a
quantum system depends on whether the coherences answer to anything, and
that is a question about the presentation --- about which equations were
withheld, and why --- rather than about the algorithm.  The framework has
made the question askable and has not answered it.

% ===================================================================
